%!TEX root =  A_WS.tex

\section{Using equations  -- Practice exercises} 

\begin{enumerate}

\item Dontrell and Kim borrowed money to buy a house on a 30-year mortgage.  
%At today's favorable interest rates, they pay  \$944 a month. (Plus taxes and insurance.)  
After $M$ months of making payments, Dontrell and Kim will still owe \$$D$ where 
$$D=\text{236,000}-\text{56,000} \ast 1.004^M$$  
$D$ is also known as the \textbf{payoff} (how much they would need to pay to settle the debt).
% Based on $j_{12}=4.8\%$

 \hfill \emph{Story also appears in 3.4 \#4}
 \begin{enumerate}
\item How much did Dontrell and Kim originally borrow to buy their house?  What value of $M$ did you evaluate at to answer the question? \vfill
\item Evaluate the equation at $M=12$ and explain what the answer means in terms of the story. \vfill
\item After making half the payments, how much money will Dontrell and Kim still owe on the house?  Will they have paid more or less (or exactly) half of the loan?  \emph{Hint:  convert 30 years into months to find the total number of payments.  Then divide by 2 to find the halfway point.} \vfill
\item The very last month they do not actually pay the regular monthly payment, just whatever balance is left on the loan.  How much will that be? \emph{Hint:  they will have made all but one of the payments.} \vfill
\end{enumerate}

\newpage %%%%%%

\item ``Rose gold'' is a mix of gold and copper.  We start with 2 grams of an alloy that is equal parts gold and copper and add $A$ grams of pure gold to lighten the color. The percentage of gold in the resulting rose gold alloy, $R$ is given by $$R = 100\left(\frac{1+A}{2+A}\right)$$
For example, if we add 0.8 grams of pure gold, then $A=.8$ and so the percentage is $$R=100\left(\frac{1+.8}{2+.8}\right)= 100 \times (1 + \underline{.8})\div(2+\underline{.8})=64.28571428\ldots \approx 64.3\%$$ 

\vspace{-.15in} %VSPACE
\hfill \emph{Story also appears in 4.1 Exercises}

\begin{enumerate}
\item Calculate the percentage of gold in the alloy if we add 1.2 grams of pure gold. \vfill \vfill
\item Fill in that and the rest of the missing values.
\begin{center}
\begin{tabular} {|c| |c |c |c |c |c |c |c  |c|}\hline
$A$ & 0 &  .4 & .8 & 1.2  & 1.6  & 2 & 3 & 4 \\  \hline
&&&&&&&& \\ 
$R$ &50.0 & \hspace{.5 in}~& 64.3 & \hspace{.5 in}~  & 72.2 & \hspace{.5 in}~& 80.0 &\hspace{.5 in}~\\
&&&&&&&& \\ \hline
\end{tabular}
\end{center}
\item Graph the function.
\begin{center}
\scalebox {.8} {\includegraphics [width = 6in] {GraphPaper.jpg}}
\end{center}
\bigskip
\item What do you think happens to the percentage of gold as we add more and more pure gold?  Try adding 10 grams, and then try adding 100 grams to check. \vfill
\end{enumerate}

\newpage %%%%%%

\item Monty hopes to grow orchids but they are fragile plants.  He will consider his greenhouse a success if at least nine of the ten orchids survive.  Assuming the orchids each survive at rate $S$, the probability his greenhouse is a success, $P$, is given by 

\begin{tabular} {ccr}
\hspace{1.55in} &$P= 10S^9-9S^{10}$ \hspace{.5in}~& \emph{Story also appears in 2.4 \#3}  \\
\end{tabular}
\begin{enumerate}
\item If the orchids are perfect ($S=1$), what is the probability of a successful greenhouse?  Explain how this answer makes sense in the story. \vfill
\item If the orchids are complete duds ($S=0$), what is the probability of a successful greenhouse?  Explain how this answer makes sense in the story. \vfill
\item Make a table comparing the probability of a successful greenhouse if the orchids each survive at rate $S = 0, .5, .8, .9, .95, $ or $1$. \vfill \vfill
\item Draw a graph of the function.
\begin{center}
\scalebox {.8} {\includegraphics [width = 6in] {GraphPaper.jpg}}
\end{center}
\end{enumerate}

\newpage %%%%%%

\item Valerie plans to do a 3-day, 50-mile walk to raise money for breast cancer research, in honor of her aunt.  Valerie's friends have pledged a total of  \$93 per mile.   
\begin{enumerate}
\item Valerie hopes to walk all 50 miles.  If so, how much money will she raise? \vfill
\item She might have to stop sooner, however. Name variables and write an equation showing how the money Valerie raises is a function of how far she is able to walk.  \vfill \vfill
\item How long will it take Valerie to walk the full 50 miles if she is able to keep a pace of 3.2 miles per hour? Write your answer in H:MM format. \vfill \vfill
\item Name the new variables and write a new equation showing how the time it takes Valerie to walk the full 50 miles depends on her pace.  \vfill \vfill
\item Good news.  Valerie walked the full 50 miles at a pace of 3.2 miles per hour.  Way to go, girl!   How much money did she raise each hour? 

 \emph{Hint: use your answers from (a) and (c) to get \$/hour.}  \vfill
\end{enumerate}

\end{enumerate}

\newpage


\noindent \textbf{When you're done \ldots}

\begin{itemize}
\item [$\Box$] Check your solutions.  Still confused?  Work with a classmate, instructor, or tutor.
\item [$\Box$] Try the \textbf{Do you know} questions.  Not sure?  Read the textbook and try again.
\item [$\Box$] Make a list of key ideas and process to remember under \textbf{Don't forget!}
\item [$\Box$] Do the textbook exercises and check your answers. Not sure if you are close enough? Compare answers with a classmate or ask your instructor or tutor.  
\item [$\Box$] Getting the wrong answers or stuck?  Re-read the section and try again.   If you are still stuck, work with a classmate or go to your instructor's office hours or tutor hours.
\item [$\Box$] It is normal to find some parts of exercises difficult, but if most of them are a struggle, meet with your instructor or advisor about possible strategies or support services.
\end{itemize}





\bigskip

\noindent \textbf{Do you know \ldots} % Using equations

\begin{enumerate} [(a)]
\item What it means to �evaluate� a function? 
\item Why some numbers are underlined in our calculation?
\item How to evaluate an function when the independent variable occurs more than once? 
\item How to generate a table or graph from an equation? 
\item What graphs of different types of functions look like? 
\item What a power, polynomial, or quadratic equation look like?
\end{enumerate}

\bigskip

\noindent \textbf{Don't forget!}


